\section{Introduction}
\label{sec:introduction}

Tax-benefit microsimulation models apply policy rules to household microdata to estimate fiscal
and distributional effects. At the national level, that usually means starting from a survey such
as the Current Population Survey (CPS), improving measured incomes and program participation, and
adjusting weights so that the resulting dataset reproduces known aggregates. Subnational work is
harder. Analysts often want estimates for states, congressional districts, or local areas, but the
underlying survey was not designed to support all of those geographies at once.

That gap matters in practice. Many reforms are proposed, administered, or debated below the
national level. A state tax proposal needs state-specific incomes, filing patterns, and benefit
participation. District-level analysis needs those same outcomes to aggregate correctly within each
district while remaining consistent with state and national totals. A usable subnational dataset
therefore has to satisfy a hierarchical calibration problem rather than a single national one.

Public survey microdata alone do not solve that problem. The CPS provides demographic structure and
program coverage, but some tax variables are missing or noisy, top incomes are imperfectly
measured, and geographic detail is limited. Administrative targets help, but they arrive from
different agencies, with different concepts, and at different geographic levels. In our setting,
the calibration problem mixes count and dollar targets across district-level units, the 50 states
plus the District of Columbia, and the nation.

The methodological goal is straightforward to state, even if the implementation is not. We want a
calibration method that can absorb a large hierarchical target system, keep weights usable for
microsimulation, and let us control how large the final dataset becomes. Exact fit is attractive,
but a production dataset also needs to remain computationally practical. A file with millions of
active cloned records is useful for detailed subnational work and still expensive to simulate. A
much smaller file is easier to ship and run, but only if the loss in fidelity is acceptable.

This paper studies an $L_0$-regularized calibration system built in PolicyEngine's US data
pipeline. The system clones CPS households across sampled geographies, builds a sparse calibration
matrix from tax-benefit simulations, and jointly optimizes positive weights and binary-like gates.
The gates give the optimizer a direct sparsity mechanism. In practical terms, that means the same
framework can target a compact national dataset or a larger subnational dataset without switching
to a different calibration procedure.

The paper's contribution is empirical as well as methodological. We do not treat generalized
regression (GREG) and iterative proportional fitting (IPF, or raking) as universal substitutes for
the full production problem. Instead, we use them as classical reference methods within scopes that
match their assumptions. The benchmark design uses shared exported calibration packages and asks
three questions. First, on a tractable benchmark, does $L_0$ perform competitively with classical
methods when all methods can plausibly run? Second, how do the methods degrade as the number of
targets increases? Third, which methods remain usable near the full production calibration problem?

The rest of the paper follows that structure. Section~\ref{sec:background} places the work in the
literature on spatial microsimulation, survey calibration, and sparse optimization.
Section~\ref{sec:data} describes the base microdata and administrative targets. Section
~\ref{sec:methodology} presents the calibration pipeline and the benchmark design. Section
~\ref{sec:results} reports the tractable benchmark, the scaling frontier, and the
production-feasibility case. Section~\ref{sec:discussion} interprets those results and discusses
their limits.
